%% Supplementary Methods: full details of main-text Sections 2.1--2.7.
%% \input by catfish_wy.tex in the Supplementary material.

\subsection*{Supplementary Methods}

These sections give the full details of Sections~2.1--2.7 of the main text, in the same order.

\subsubsection*{SM1 Gene-level inputs and component tests}

For a pathway $S$ containing $m=|S|$ genes, let $p_g$ denote the gene-level
association $p$-value for gene $g\in S$ and $Z_g$ the corresponding gene-level $Z$-statistic. We write $\mathbf{p}_S=(p_1,\dots,p_m)$ and $\mathbf{Z}_S=(Z_1,\dots,Z_m)$. Gene-level statistics are obtained with MAGMA \citep{deLeeuw2015}, which aggregates SNP-level GWAS evidence across variants assigned to a gene while accounting for LD, gene size, and local marker density, and which additionally returns the gene--gene correlation matrix induced by LD. These quantities are the common inputs for every CATFISH component test and for both calibration layers.

Each pathway is assessed using five complementary component tests (Table~\ref{tab:components}): ACAT \citep{Liu2019ACATV,Liu2020ACAT}, Fisher \citep{Fisher1932}, adaptive soft-thresholding TFisher \citep{Zhang2020TFisher}, minP, and Stouffer \citep{Stouffer1949}. Together, they cover the signal archetypes of Section~3.1 and apply directly to gene-level $p$-values or $Z$-scores without external weights or pathway-specific tuning. We combine the component $p$-values into a single omnibus statistic via an ACAT-type cross-method operator $\mathcal{O}(\cdot)$. This omnibus is a pathway-level summary of evidence across components, not an additional gene-level test.

Four components (ACAT, Fisher, minP, Stouffer) are fully determined by gene-level $p$-values and require no tuning. Soft TFisher is the exception: it adapts over a grid of truncation thresholds $\tau$, which determines which genes contribute. This is the only tunable choice in CATFISH, and we justify it below.

Following \citet{Zhang2020TFisher}, the grid spans small, intermediate and large thresholds rather than concentrating at one end: we use \(\tau\in\{10^{-2},\,5\times10^{-2},\,5\times10^{-1}\}\) and take the minimum calibrated TFisher \(p\)-value over the grid. The original authors' grid also includes \(\tau=1\), which we omit because \(\tau=1\) is Fisher's method and Fisher is already carried as a separate component. Spanning the range matters because the number of genes surviving truncation is \(\mathrm{Binomial}(m,\tau)\): a grid confined to small \(\tau\) leaves \(m\tau\ll1\) at realistic set sizes, the statistic is then zero on a large fraction of pathways, and the component is degenerate precisely where it is most often applied.

A threshold is informative mainly when \(m\tau\) is not far below one, so which \(\tau\) values are active depends on set size: \(\tau=5\times10^{-1}\) is informative from the smallest sets upward, whereas \(\tau=10^{-2}\) requires \(m\approx100\). Adapting across \(\tau\) rather than fixing a single value is what allows one component to serve both sparse and diffuse architectures, and users can assess the grid for their own database using \(m\tau\).

\smallskip\noindent\textit{Undefined components.} When no gene in a pathway falls below any threshold in the grid, the soft TFisher statistic is zero at every \(\tau\) and its one-sided \(p\)-value is exactly one. We treat this as \emph{undefined} rather than as evidence against enrichment, and the omnibus for that pathway is formed from the components that are defined. The distinction is not cosmetic. The ACAT operator maps each component \(p\)-value through \(\tan\{(1/2-p)\pi\}\), so a value at the upper boundary contributes a large negative term; a component that could not be computed would otherwise vote against the pathway and bias the omnibus toward the null. The rule is applied per pathway rather than by set size, because a large pathway in which no gene is individually strong is equally uncomputable. It does condition the component on having at least one gene below a threshold, so where TFisher does contribute its \(p\)-values are stochastically small; because the resampling null is generated under the identical rule, that conditioning appears in the null distribution and is absorbed by the calibration. An analytic omnibus \(p\)-value under this rule would be anti-conservative, whereas the resampled one is not.

\smallskip\noindent\textit{Precision limit.} The soft TFisher implementation returns the null distribution function of the soft-thresholding statistic, so the one-sided \(p\)-value is formed as one minus that value. In double precision the difference is exactly zero once the distribution function exceeds \(1-10^{-16}\), and a strongly enriched pathway would then report \(p=0\). We cap the soft TFisher \(p\)-value at \(10^{-16}\), the smallest value the computation can resolve, rather than at the general lower floor applied to components whose tails are computed in closed form. A cap below the precision limit asserts resolution the statistic does not have; Section~4.6 shows what that does to the family-wise threshold.

\subsubsection*{SM2 Two sources of dependence}

CATFISH statistics are correlated at two levels. The two arise from different
causes, and both have to be addressed before any pathway $p$-value means what
it claims to.

\emph{Within a pathway}, component tests are dependent for two reasons: each is a deterministic function of the same $\mathbf{p}_S$, inducing cross-method dependence even if genes were independent; and gene-level evidence is correlated via LD, overlapping SNP-to-gene assignments, and shared local genomic architecture. Thus, any analytic reference distribution assuming independence is wrong for every component and typically anti-conservative, because correlated genes contain less information than the same number of independent genes.

\emph{Between pathways}, pathway-level $p$-values are correlated because databases assign genes to multiple sets. A gene shared across pathways propagates its (possibly extreme) null draw to all of them, shifting their evidence in the same direction. The strength and pattern of this dependence depend on database structure: hierarchies with nested sets create widespread but weak dependence, whereas largely disjoint curated collections create rare but strong dependence among overlapping sets. Consequently, independence-based multiple-testing corrections (Bonferroni or Benjamini--Hochberg under its positive-dependence conditions) need not achieve nominal error rates; the discrepancy is driven by the database, not the method.

These are often treated as separate problems, but in CATFISH both reduce to the same object: a vector of correlated gene-level statistics. A single genome-wide null draw, carrying the estimated LD, is scored through all pathways at once. Within a pathway, it provides null components and accounts for within-pathway dependence; across pathways, it provides a null for the family minimum and accounts for between-pathway dependence. No correlation structure is estimated or modeled; both are inherited. Section~2.3 develops this.

The construction relies on three assumptions: (i) \emph{Approximate normality}: under the null, the pathway-level vector satisfies $\mathbf{Z}_S\sim\mathcal{N}(0,R_S)$. (ii) \emph{Correlation-matrix adequacy}: correlation matrices from MAGMA-reported gene pairs provide an LD-informed working approximation to the unobserved gene-level covariance; unreported pairs are set to zero, and adequacy is evaluated via calibration and robustness analyses below. (iii) \emph{No residual structural confounding}: gene-level inputs are adjusted so pathway signals are not primarily driven by gene size or SNP density.

\subsubsection*{SM3 Calibration and family-wise control from a shared correlated null}

CATFISH obtains two different $p$-values for every pathway: a marginal one,
which asks whether that pathway is enriched, and a family-wise one, which asks whether evidence this strong would appear anywhere among the $K$ pathways tested by chance. Both are read from a single set of resampled null draws. The procedure has four steps.

\paragraph{Step 1: generate the null.}
Gene--gene correlations are read from MAGMA's \texttt{.genes.raw} output as a list of gene pairs and their correlations. For each chromosome $c$ containing $m_c$ genes we form an $m_c\times m_c$ correlation matrix $R_c$ by starting from the identity and filling in the reported pairs,
\[
R_c = I_{m_c} + \sum_{(i,j)\in\mathcal{P}_c}
      r_{ij}\left(e_ie_j^{\top}+e_je_i^{\top}\right),
\]
where $\mathcal{P}_c$ is the set of reported pairs for chromosome $c$; pairs not in $\mathcal{P}_c$ remain zero. We compute correlations only for gene pairs close enough to share LD, so $R_c$ is banded by construction: zeros for distant genes are approximations, not imputations, whereas zeros for nearby but unreported pairs introduce approximation error (see Section~5). Genes on different chromosomes are treated as independent, which is structurally correct. If $R_c$ is not positive definite, we floor eigenvalues at $10^{-6}$, reconstruct, and renormalise the diagonal; the implementation reports affected chromosomes.

Writing $R_c=U_c^{\top}U_c$ for the Cholesky factor, replicate $b$ draws
\[
\mathbf{Z}^{(b)}_c = U_c^{\top}\boldsymbol{\varepsilon},\qquad
\boldsymbol{\varepsilon}\sim\mathcal{N}(0,I_{m_c}),
\]
independently per chromosome, and the concatenation $\mathbf{Z}^{(b)}$ is one genome-wide null realisation carrying the estimated LD structure. Gene-level null $p$-values follow by the Gaussian copula $p_g^{(b)}=1-\Phi(Z_g^{(b)})$, matching MAGMA's one-sided sign convention. This is repeated for $b=1,\dots,B$.

We draw genome-wide, not pathway-by-pathway, because Step~4 requires all $K$ pathways to be scored on the \emph{same} realisation so shared genes perturb them jointly. Drawing each pathway from its own correlation submatrix would be marginally correct but would remove the dependence structure Step~4 needs.

\paragraph{Step 2: score every pathway on every replicate.}
For each replicate we recompute all five component statistics and the omnibus operator $\mathcal{O}(\cdot)$ for every pathway in the family $\mathcal{F}=\{S_1,\dots,S_K\}$, giving a $B\times K$ matrix of null omnibus values
\[
Q^{(b)}_k = \mathcal{O}\!\left(\mathbf{p}^{(b)}(S_k)\right),
\qquad b=1,\dots,B,\quad k=1,\dots,K .
\]
Because all components are recomputed from the same draw, $Q$ preserves within-pathway gene dependence, cross-method dependence, and—since all $K$ pathways share the draw—between-pathway dependence. The observed values $q^{\mathrm{obs}}_k=\mathcal{O}(\mathbf{p}^{\mathrm{obs}}(S_k))$ are computed the same way from the real data. Applying $\mathcal{O}$ to the observed components without resampling yields the uncalibrated analytic omnibus $\mathrm{Omnibus}_{\mathrm{An}}(S_k)=q^{\mathrm{obs}}_k$, which assumes gene independence and serves as the baseline.

\paragraph{Step 3: the marginal \texorpdfstring{$p$}{p}-value.}
Compare each pathway's observed statistic with its own null values across the $B$ replicates. The LD-aware marginal $p$-value is the Monte Carlo tail probability
\[
\tilde p_k \;=\;
\frac{1+\sum_{b=1}^{B}\mathbf{1}\!\left(Q^{(b)}_k\le q^{\mathrm{obs}}_k\right)}{B+1}.
\]
This quantity is computed from the genome-wide draw rather than from the pathway-specific submatrix $R_{S_k}$; because $R_{S_k}$ is simply the $R_c$ submatrix for that pathway's genes, both use the same correlation information, so nothing is lost. We call $\tilde p_k$ the per-pathway calibrated $p$-value and use it to rank pathways. It corresponds to what is often called MVN calibration; we avoid that label because the multivariate normal model is the generator of the null draws, not a separate calibration step, and the same draws serve Step~4.

\paragraph{Step 4: the family-wise \texorpdfstring{$p$}{p}-value.}
Within each replicate, record the smallest statistic across the whole family,
\[
M^{(b)} = \min_{1\le k\le K} Q^{(b)}_k ,
\]
and compare each observed statistic with the resulting distribution. This gives the Westfall--Young adjusted $p$-value \citep{Westfall1993}
\[
\tilde p^{\mathrm{WY}}_k \;=\;
\frac{1+\sum_{b=1}^{B}\mathbf{1}\!\left(M^{(b)}\le q^{\mathrm{obs}}_k\right)}{B+1},
\]
Rejecting all pathways with $\tilde p^{\mathrm{WY}}_k\le\alpha$ controls the family-wise error rate at level $\alpha$ under the complete null, whatever the dependence among pathways, because a false rejection anywhere in the family requires the family minimum to cross the threshold. Three further points apply. First, the control is in the strong sense, not only under the complete null. Westfall--Young procedures control the family-wise error rate strongly when the joint distribution of the statistics for the true null hypotheses does not depend on which other hypotheses are false---the condition of subset pivotality \citep{Westfall1993}. Two features of the construction give it. Each pathway statistic is a deterministic function of its own genes' $z$-scores only; and under the model of Step~1, with association represented as a shift in the means of the associated genes, the genes of a pathway containing no associated gene are jointly $\mathcal{N}(0,R)$, restricted to those genes, whatever happens elsewhere. The observed statistics of the null pathways therefore have exactly their complete-null joint distribution, and because the minimum over any subset of pathways is never smaller than the minimum over the whole family, the probability that any null pathway crosses the threshold is at most $\alpha$. The argument requires the soft TFisher truncation grid to be fixed in advance: a grid chosen from the empirical distribution of all gene-level $p$-values, as in data-adaptive implementations, would make every pathway's statistic depend on genes outside it. Second, the null is parametric: it is generated from the multivariate normal model of Step~1 rather than by permuting phenotypes, so the procedure is Westfall--Young in construction but inherits the adequacy of that model. Third, the counting form of the adjusted $p$-value above, in which null minima equal to the observed value count as at least as extreme, is required rather than a matter of taste. A threshold taken as a quantile of $M$ and applied with an inclusive comparison rejects every observation tied at that quantile, and if the omnibus loses resolution in the tail so that many null minima coincide, the empirical family-wise error rate can approach one (Section~4.6).

Between-pathway correlation is never explicitly estimated or modelled. Because all $K$ pathways are evaluated on the same draw in Step~2, genes shared across pathways perturb them jointly, as in the observed data, and this dependence is reflected in the distribution of $M$. The threshold $\hat t_\alpha=\mathrm{quantile}_\alpha(M)$ is the family-wise significance cutoff implied by the pathway database in use; its ratio to the Bonferroni cutoff $\alpha/K$ quantifies the database’s departure from independence at no extra cost, provided it is read on calibrated statistics; on the analytic scale the ratio also absorbs any miscalibration of the per-pathway statistic (Section~4.1).

Steps~3 and~4 use the same matrix $Q$ and require no additional resampling. Their outputs relate as a raw and adjusted $p$-value: $\tilde p_k$ assesses pathway $S_k$ alone, while $\tilde p^{\mathrm{WY}}_k$ accounts for testing $K$ pathways. Both are reported; neither replaces the other.

\paragraph{Size standardisation.}
A minimum across tests is meaningful only when the tests share a scale. The raw omnibus does not: Fisher’s statistic increases with $m$, so the Step~4 family minimum would mechanically favor the largest sets. We therefore also report a variant that first standardizes each observed value within its own column of $Q$ as a $z$-score, $z_k=(q^{\mathrm{obs}}_k-\mu_k)/\sigma_k$, where $\mu_k$ and $\sigma_k$ are that column’s mean and standard deviation, and then takes the row minimum. This rescales values for the minimum operation; it is not a pathway-size bias correction, and Section~4.1 shows it does not behave as one. In the simulations of Section~2.6 every pathway in a configuration has the same size and dependence structure, so the pathways are exchangeable, share one marginal null, and the row minimum is taken on the raw scale; standardisation is needed only when pathway sizes differ, as in real data. We standardise by location and scale rather than by each pathway's empirical null distribution function. The latter would put pathways on an exactly common scale, but would resolve each to only $1/(B+1)$, which is coarser than the family minimum requires and would reintroduce ties.

\paragraph{Competitive test.}
The Westfall--Young test above is self-contained: its null hypothesis is that no
gene in the pathway is associated with the trait. CATFISH also tests the
competitive hypothesis that the genes of a pathway are no more associated than
genes elsewhere in the same GWAS. The omnibus statistic and both readings of the
null matrix (each pathway's calibrated $p$-value and the family-wise
$p$-value from the minimum of the standardised statistics) are unchanged; only
the null replicates differ. Genes are placed in genome order (chromosome, then
start position) and joined end to end. In each replicate every gene takes the
observed score of the gene $s$ places further along this order, wrapping at the
end, for one random shift $s$ shared by all pathways, and every pathway is
rescored. A pathway's genes therefore keep their spacing, so genes in linkage
disequilibrium with one another move together, while their scores come from a
random position in the genome; the genome-wide distribution of gene scores,
including any polygenic signal, is present in the null as well as in the
observed statistic. Because one shift is applied to all pathways, overlap
between pathways is preserved within each replicate and the minimum over
pathways gives strong family-wise control. Shifts are drawn without replacement,
so at most $G-1$ replicates exist for $G$ genes. The null does not match genes
on size or SNP count. A pathway is called by both tests when
$\max(p_{\mathrm{FWER}}^{\mathrm{self}}, p_{\mathrm{FWER}}^{\mathrm{comp}})\le\alpha$,
an intersection--union test that controls the family-wise error of the joint
claim at $\alpha$ without further correction \citep{Berger1982}.

\subsubsection*{SM4 Numerical implementation in log space}

The omnibus statistic must be computed entirely in log space. This is a correctness requirement, not an optimisation, and it prevents silent failure.

The smallest positive normalized double is $\approx2.2\times10^{-308}$. In a GWAS with $n\approx1.3$ million, gene-level $p<10^{-300}$ is common: in the GLGC LDL-cholesterol data (Section~3.4), the published $p$-value is exactly zero for 473 chromosome~1 variants, including the \textit{SORT1} lead rs12740374 at $-\log_{10}p=1725$. Combining even a few such genes with Fisher’s method yields a pathway statistic whose double-precision $p$-value underflows to zero. All affected pathways then share the same floor value, ties are broken by storage order, and the ranking the analysis aims to produce is silently corrupted. In our first implementation, the top-ranked pathways were (on inspection) alphabetical.

Five changes are needed. (i) Compute all component statistics on $\ell_g=\log p_g$ and return log $p$-values via \texttt{stats::pchisq(\dots,\,log.p=TRUE)} and \texttt{stats::qnorm(\dots,\,log.p=TRUE)}. (ii) Evaluate the ACAT sum using the small-$p$ asymptote: for $p<10^{-15}$, use $\tan\{\pi(\tfrac12-p)\}\approx 1/(p\pi)$, accumulating those terms as a log-sum-exp of $-\ell_g-\log\pi$; use the trigonometric form for the rest. (iii) Replace $p=\tfrac12-\arctan(T)/\pi$ with $p\approx1/(T\pi)$ for $T>10^{6}$, where cancellation near $\tfrac12$ dominates truncation error. (iv) Since ACAT can be negative when most inputs exceed $0.5$, handle that case in a separate branch computing $\tfrac12-\arctan(T)/\pi$ directly; clamping would set $p=0.5$ and destroy ordering in the upper half. (v) The minP inversion $1-(1-q)^m$ underflows for $q\lesssim10^{-16}/m$, causing minP to drop out; use the $m q$ asymptote whenever $\log m+\ell_{\min}<\log 10^{-8}$.

These changes let the pipeline represent pathway $p$-values down to $e^{-745}\approx10^{-324}$ and, crucially, preserve their \emph{ordering} well below where probability-scale computations fail. They are implemented in the scripts used for the real-data analyses. The simulations of Section~2.6 instead call the component functions of the CATFISH package directly, so that they validate the code a user runs; those functions work on the probability scale with a lower floor of $10^{-300}$, which is adequate there because no simulated null approaches it. Soft TFisher is subject to its own, much higher, precision limit (Section~2.1).

\subsubsection*{SM5 Candidate-gene prioritization}

The output a biologist acts on is a short list of genes, so CATFISH ranks genes by the sum of three evidence layers:
\[
\mathrm{Score}(g)=S_{\mathrm{GWAS}}(g)+S_{\mathrm{MAGMA}}(g)+S_{\mathrm{PATH}}(g).
\]
Each layer enters by rank rather than by its raw statistic,
$S_\bullet(g)=-\log_{10}\!\left(r_\bullet(g)/N_\bullet\right)$, where
$r_{\mathrm{GWAS}}$ is the gene's rank by minimum SNP-level $p$-value,
$r_{\mathrm{MAGMA}}$ its rank by gene-level $p$-value, and $r_{\mathrm{PATH}}$ the rank of its best pathway by the size-standardised Westfall--Young statistic $z_k$ of Section~2.3; rank one is the strongest evidence. We rank pathways on $z_k$ for two reasons. It is the statistic the family-wise test is computed on, so the gene score uses the same pathway ordering that is reported; and it measures each pathway against its own correlated null, whereas the analytic omnibus assumes independent genes and overstates pathways whose genes are in LD (Section~3.2). The family-wise $p$-value gives the same ordering but is floored at $1/(B+1)$, which ties the strongest pathways. Genes missing a layer score zero for it, so an absent layer neither helps nor penalises.

We use ranks instead of $-\log_{10}p$ to make layers with different tail behaviour comparable and to avoid the $p$-value floors in Section~2.4. Because $r_\bullet(g)/N_\bullet$ is the fraction of genes ranked at least as high as $g$ on a layer, each term is $-\log_{10}$ of a rank-based empirical $p$-value, and the unweighted sum is Fisher's combination of the three layers. Equal weights need no known causal genes, so the score applies unchanged to species and pathway databases without curated gene sets. An earlier version of the score weighted the pathway layer by $0.5$ and added a hub term $0.3\log_{10}(n_{\mathrm{path}}(g)+1)$ rewarding membership in several pathways; we dropped the hub term because it is not on the rank scale of the other layers and it rewards well-annotated genes, which is the main route by which annotation bias could enter the score. We assessed the pathway and hub weights over a grid, $w_{\mathrm{path}}\in\{0,0.25,0.5,0.75,1,1.5,2\}$ and $w_{\mathrm{hub}}\in\{0,0.1,\dots,0.5\}$, in a truth-planted simulation and in the real analyses of Section~3.5 (Supplementary Figure~\ref{fig:weights}). The simulation draws $400$ replicate datasets of $500$ genes and $50$ pathways, plants causal genes concentrated in five pathways, and generates correlated GWAS-, MAGMA- and pathway-level evidence for them \citep{Tandukar2025Thesis}. Rankings were stable across the grid: the Spearman correlation of each full ranking with the equal-weight ranking was at least $0.98$. Recovery of planted causal genes in the top twenty rose with the pathway weight, from $0.71$ without the pathway layer to $0.83$ with equal weights and $0.88$ at twice the pathway weight, and the hub term added about $0.01$.

\emph{Where the gain comes from.} Two checks separate the contribution of the
pathway layer from that of the association layers. First, genes are ranked with
the pathway weight set to zero, so the score reduces to the GWAS and MAGMA
layers. Second, a gene's own association contributes to the evidence of every
pathway that contains it, so a gene could be raised by its own signal passing
through its pathway. The pathway evidence of each gene being checked (for
example, the positive controls of a trait) is therefore recomputed with that
gene left out: for each tested pathway $P$
containing gene $g$, $P$ and $P$ without $g$ are scored in one Westfall--Young
run on the same null draws, the change in $z_k$ is added to $P$'s statistic from
the full analysis, the strongest of these becomes $g$'s pathway evidence, and
its pathway rank, score and overall rank are recomputed with every other gene
unchanged (\texttt{wy\_logo.R}, \texttt{39\_logo\_recovery.R}). Reduced sets
are allowed one gene fewer than the minimum pathway size, so that removing a
gene does not drop its pathway. If the gene keeps its rank, its lift comes from
the other genes of its pathways; if it falls back to its rank without the
pathway layer, the lift was its own.

The score is a heuristic ranking metric, not an inferential model: it has no null distribution or $p$-value. We also report per-layer ranks, because the GWAS and MAGMA layers use the same variants and are not independent evidence.

\subsubsection*{SM6 Simulation design}

Simulations are run on the omnibus statistic under its null, except where a signal is injected deliberately (item~6). Two dependence designs are used. In the \emph{shared-universe} design, pathways are overlapping subsets of one gene universe: $G=3000$ genes in independent blocks of 50 with exchangeable within-block correlation $\rho$, standing in for chromosomes and LD blocks, from the first 600 of which we draw $K=150$ pathways of $m$ genes each at random. Each gene therefore belongs to several pathways, a pathway spans several blocks, and a single genome-wide draw moves every pathway containing a given gene---the between-pathway dependence the family-wise procedure is designed to control. In the \emph{exchangeable} design, used only in item~1, each pathway is drawn independently and all $m$ of its genes are mutually correlated at $\rho$. That design has no between-pathway dependence; it isolates within-pathway correlation at an intensity well beyond what block-structured LD produces, and serves as a stress test. Each study reports two summaries of the null matrix: per-pathway marginal type~I error (by columns) and family-wise error rate (by rows), with Bonferroni as the family-wise comparator. Thresholds are fitted on one set of draws and evaluated on a held-out set, so rates are not correct by construction. Component statistics are computed with the CATFISH package functions, using the soft TFisher grid, undefined-component rule and precision cap of Section~2.1. We do not simulate genotypes or phenotypes or run a full GWAS; that is future work (Section~5) and is not needed for the questions here. Each study corresponds to a script in \texttt{catfish2/scripts/} and writes a CSV to \texttt{catfish2/results/} from which the tables and figures are generated.

\begin{enumerate}
\item \textbf{Null calibration and type~I error} (\texttt{01\_null\_calibration.R}), exchangeable design. Under the global null we compare four decision rules: the analytic independence-based $p$-value; per-pathway calibration against a null bank shared by all pathways of a given size; Bonferroni; and the Westfall--Young adjusted $p$-value of Step~4, with the null distribution of the family minimum estimated on one half of the simulated families and evaluated on the other. We vary $\rho\in\{0,0.3,0.6\}$ and $m\in\{5,20,50\}$ and report type~I error at $\alpha\in\{0.05,0.01\}$, family-wise error, the genomic-control inflation factor $\lambda=\mathrm{median}(\chi^2_1\text{-transformed }p)/0.4549$, and the number of distinct values taken by the family minimum, which detects loss of resolution directly.

\item \textbf{Component-breakage stress test} (\texttt{02\_breakage\_stress.R}).
We miscalibrate $k$ of five components by replacing each affected component’s $p$-value with $p^{\gamma}$, and test whether the omnibus and its calibration remain valid. We run both corruption directions: for $0<p<1$, $\gamma<1$ yields $p^{\gamma}>p$ (conservative) and $\gamma>1$ yields $p^{\gamma}<p$ (anti-conservative), the latter being most relevant for false positives. We use $\gamma\in\{0.7,1.4\}$; sweep $k=0,\dots,4$; evaluate up to three component combinations per $k$; and cross with the same LD regimes at $m\in\{20,50\}$. We apply corruption to both null replicates and observed statistics, allowing resampling calibration to absorb it.

\item \textbf{Leave-one-out component sensitivity} (\texttt{03\_adaptive\_loo.R}).
We remove each component in turn, recompute the null bank and the calibration from scratch, and measure the change in type~I error and $\lambda$. This isolates how much each component contributes to, or destabilizes, the calibrated omnibus, and identifies components whose removal would change inference.

\item \textbf{Per-component calibration} (\texttt{09\_component\_calibration.R}).
Rather than combining, we calibrate each of the five components against its own null and report analytic and calibrated type~I error, family-wise rate and $\lambda$ per component. This identifies which statistics LD actually breaks and whether a single calibrated component is usable in place of the omnibus when the signal architecture is known in advance.

\item \textbf{Robustness to an incomplete correlation file}
(\texttt{08\_missing\_correlation.R}). Data are simulated under the true correlation, but calibration (null bank, marginal $p$-values, and Westfall--Young threshold) uses a matrix where a fraction $\{0,0.25,0.5,0.75,1\}$ of off-diagonal entries is set to zero, mimicking unreported gene pairs. Evaluation draws are shared across missingness levels; only calibration varies. For real data, \texttt{10\_pair\_coverage.R} places a correlation file on this scale by reporting the fraction of within-pathway gene pairs on the same chromosome and close enough to be correlated for which MAGMA reports a value.

\item \textbf{Archetype power} (\texttt{17\_archetype\_power\_wy.R}), shared-universe design with $\rho=0.3$. One pathway in each family carries signal shaped as one of the five archetypes of Section~3.1 (Table~\ref{tab:archetype_summary}); the remaining $K-1$ are null. Signal is a mean shift on the gene-level $z$ scale applied to a fixed fraction of the pathway's genes, at positions permuted in every replicate: SSS shifts $15\%$ of genes by $3.0$; MDS $60\%$ by $1.1$; WDE $90\%$ by $0.55$; MSA $10\%$ by $1.8$ and a further $60\%$ by $1.1$; and SGP $5\%$ by $4.0$, with at least one gene in every layer. Shifts are multiplied by an effect scale over a grid running from undetectable to saturated, at $m\in\{5,20,50\}$. For each component and for the omnibus we record the probability that the signal pathway clears its own Westfall--Young threshold. Regret is the difference in that probability between the best single component and the omnibus in the same cell.
\end{enumerate}

Unless noted, each configuration uses $B$ null draws for the calibration bank and separate replicate datasets for evaluation; both settings are defined at the top of each script, allowing a reduced smoke test via \texttt{bash run\_all.sh quick}.

\paragraph{Simulation with known truth on real linkage disequilibrium.}
We compared CATFISH with MAGMA gene-set analysis on GWAS summary statistics simulated under real LD (\texttt{scripts/hpc/simulation\_check\_catfish}), ensuring both methods used identical SNP statistics and MAGMA gene scores. Genotypes from a reference panel of $n$ individuals were restricted to selected chromosomes/SNPs and standardised to $X$ ($n\times M$). With $u=e/\sqrt{n}$ and $e\sim\mathcal{N}(0,I_n)$, SNP statistics $z=X^\top u$ had covariance equal to the panel correlation matrix. 

We added signal at causal SNPs by assigning each causal SNP $k$ an expected statistic $\lambda_k$ (random sign) and propagating it to SNP $j$ within 1~Mb as $r_{jk}\lambda_k$, where $r_{jk}$ is the panel correlation; limiting propagation to local LD avoids inflation from chance long-range correlations in a finite panel. Each causal gene was assigned one causal SNP, chosen uniformly from its mapped SNPs. 

We simulated two nulls: (i) no causal SNPs; and (ii) a polygenic background where a fixed fraction of genes had one modest-effect causal SNP, with no pathway enrichment. Other scenarios added signal in one target pathway on top of this background, following the five archetypes in Table~\ref{tab:archetype_summary}: a specified fraction of pathway genes was made causal, with $|\lambda|=\delta$ (archetype effect) times a shared effect scale. Background genes were sampled outside the target pathway. Target pathways (three prespecified sizes) were chosen from pathways unrelated to the mimicked trait. Each replicate was analysed with MAGMA gene analysis (SNP-wise mean), CATFISH (self-contained, competitive, and both), and MAGMA gene-set analysis (competitive and self-contained via \texttt{--model self-contained}). CATFISH used Westfall--Young FWER control at $p\le0.05$; MAGMA used Bonferroni, $p\le0.05/K$ over $K$ tested pathways. Power was the fraction of replicates in which the target pathway was called. An error was calling a pathway with no genes within 1~Mb of any target causal SNP; with no target pathway, every call was an error. Family-wise error was the fraction of replicates with at least one error. Under a polygenic background, self-contained tests are expected to call pathways containing truly associated background genes, so errors were assessed for competitive tests and the joint call only.

Simulations used the 1000 Genomes EUR lipid panel (503 individuals), chromosomes 1–10, and HapMap3 SNPs with MAF $\ge0.01$ (748,822 SNPs). Genes were annotated with MAGMA ±10 kb windows (9,645 genes); Reactome was restricted to these genes, yielding K=1,021 pathways (10–200 genes). The background set 2% of genes as causal with |λ|=3. An effect scale of 3 was chosen from a pilot (2, 4, 6) to avoid no detection or saturation. Target pathways contained 13, 36, and 97 genes. The 17 scenario–size combinations had 50 replicates each (850 total). CATFISH used B=2,000 null replicates per test, five components with τ∈{0.01,0.05,0.5}, and gene-level genomic control.

\subsubsection*{SM7 Real-data analyses}

\textbf{Maize grain nitrogen.} Gene-level statistics were computed with MAGMA from mixed linear model (MLM) GWAS results for grain nitrogen content, assigning SNPs to genes within $\pm 10\,\mathrm{kb}$. Pathways were from CornCyc \citep{Hawkins2021CornCyc,Caspi2008BioCyc}, yielding $172$ pathways after size filtering. Gene--gene correlations were extracted per chromosome from MAGMA \texttt{.genes.raw}; we confirmed all ten chromosomes were included after an earlier parsing error captured only chromosome~1 and treated $85\%$ of the genome as independent. Candidate genes were scored with the three-layer method (Section~2.5) across all $34{,}873$ annotated genes.

\textbf{Human LDL and HDL cholesterol.} We used the Global Lipids Genetics
Consortium LDL- and HDL-cholesterol summary statistics from the European-ancestry meta-analysis ($n\approx1.3$ million) \citep{Graham2021GLGC}. Both traits were processed identically, as described below, and the ancestry-specific meta-analyses of the same consortium (Section~3.4) were processed in the same way.

\emph{Variant filtering and the $p$-value floor.} Variants were restricted to rsIDs in the 1000 Genomes European LD reference panel. The released \texttt{pvalue} column underflows to zero at top loci, so $p$-values were reconstructed from \texttt{pvalue\_neg\_log10}. For MAGMA, which requires probabilities via \texttt{--pval}, values were floored at $10^{-300}$; this does not affect MAGMA because the floor is well below its gene-level resolution. The candidate-gene GWAS layer is affected because it ranks genes by their minimum variant $p$-value, and flooring creates ties among the strongest genes. It therefore uses an unfloored $-\log_{10}p$ file produced by \texttt{11\_munge\_neglog10.sh}.

\emph{Gene analysis.} Variants were mapped to genes using NCBI37.3 locations with a $10\,\mathrm{kb}$ upstream/downstream window. MAGMA was run per chromosome with \texttt{--batch} and combined using \texttt{--merge} to retain gene--gene correlations in the merged \texttt{.genes.raw} file; concatenating per-chromosome outputs would lose these correlations.

\emph{Pathways.} Pathway inference used Reactome from MSigDB C2:CP ($1{,}390$ sets after size filtering); HumanCyc was also converted to CATFISH format for comparison. We report Reactome because its hierarchical structure induces strong between-pathway dependence, the setting targeted by the family-wise procedure.

\emph{Candidate genes.} The GWAS layer uses MAGMA's own \texttt{.genes.annot} file rather than re-deriving a window, ensuring the GWAS and gene layers use identical variant sets. The pathway layer is ranked by the size-standardised statistic $z_k$ instead of the standardised Westfall--Young $p$-value. Both produce the same ordering because the $p$-value is a monotone transform of $z_k$, but the $p$-value is floored at $1/(B+1)$, tying the top results at the $B$ used here.

\emph{Positive controls.} Lipid candidate genes were evaluated against the Open Targets gold-standard genes of each trait \citep{Mountjoy2021L2G} from the 2019 and 2023 releases (LDL-C: LDL cholesterol measurement and hypercholesterolaemia; HDL-C: HDL cholesterol measurement), restricted to genes scored by both CATFISH and PoPS: $15$ for LDL-C and $18$ for HDL-C (\texttt{45\_benchmark\_topk\_all.R}, \texttt{45\_gold\_genes.csv}). They were curated independently of this study, are used only for assessment, and the score does not use them.

\emph{Comparison with MAGMA.} We compared CATFISH with MAGMA v1.10 \citep{deLeeuw2015} using the same merged \texttt{.genes.raw} file for each lipid set and the same Reactome \texttt{.sets.annot} file. MAGMA’s default competitive test regresses gene $z$-scores on set membership, adjusting for gene size, gene density, sample size, and minor allele count, to assess whether set genes are more associated than other genes. With \texttt{--model self-contained}, MAGMA instead reports a self-contained test matching CATFISH’s question. We restricted results to the $1{,}390$ pathways tested by CATFISH and applied Bonferroni correction across these sets, since MAGMA provides no family-wise procedure. MAGMA uses its internally computed gene statistics without the centring described below.

\textbf{Sorghum dry biomass.} We analysed dry biomass yield (dry tons per acre) in the sorghum Bioenergy Association Panel (BAP) \citep{Brenton2016}: $327$ lines with measurements. Variant--gene LD was computed from the panel genotypes ($363$ lines; minor allele frequency $\ge0.01$). Associations were tested using GEMMA’s univariate linear mixed model \citep{Zhou2012GEMMA}, yielding $\nSorSnps$ variants passing MAGMA’s sample-size filter.

\emph{Gene analysis.} We ran MAGMA v1.10 \citep{deLeeuw2015} separately per chromosome, assigning variants to BTx623 v3 gene models with a $25\,\mathrm{kb}$ upstream/downstream window. Because MAGMA computes gene--gene correlations only for nearby genes on the same chromosome, the ten per-chromosome \texttt{.genes.raw} files together capture the full correlation structure. We used MAGMA’s SNP-wise mean model (as for the lipid traits); SNP-wise top and multi-model $p$-values were inflated with gene size and were excluded (Supplementary Figure~\ref{fig:sorghum_pcol}).

\emph{Pathways and calibration.} Pathways came from SorghumBicolorCyc \citep{Hawkins2021CornCyc}, limited to $5$--$200$ genes ($\nSorDryNpw$ pathways). Westfall--Young calibration matched the other analyses ($B=10{,}000$; same truncation grid and seed).

\emph{Candidate genes.} The GWAS layer uses the variant $p$-values analysed by MAGMA and assigns them to genes using MAGMA’s gene windows, ensuring the GWAS and gene layers use identical variant sets (as for lipid traits).

\emph{Positive controls.} We pre-specified eight cloned genes affecting flowering time and plant height (key biomass determinants in this photoperiod-sensitive panel): \textit{Ma1}/\textit{PRR37} \citep{Murphy2011PRR37}, \textit{Ma3}/\textit{PHYB} \citep{Childs1997PHYB}, \textit{Ma5}/\textit{PHYC}, %% TODO: add the Ma5/PHYC citation
\textit{Ma6}/\textit{GHD7} \citep{Murphy2014Ghd7}, \textit{Ma2} \citep{Casto2019Ma2}, \textit{Dw1} \citep{Hilley2016Dw1}, \textit{Dw2} \citep{Hilley2017Dw2}, and \textit{Dw3} \citep{Multani2003Dw3}. Because SorghumBicolorCyc focuses on metabolism, none appear in the tested pathways; they assess only the GWAS and MAGMA layers.

\emph{Phenotype-permutation null.} Section~2.6 simulations draw gene $z$-scores from the same correlation model as the null, testing the procedure under its assumed model. To also test the model, we permuted dry biomass among the $327$ lines $100$ times. Permutation breaks any association while preserving genotypes, LD, kinship, and the phenotype distribution. Each permuted phenotype was run through the full analysis pipeline: GEMMA mixed model, MAGMA per chromosome with the same windows and SNP-wise mean model, and Westfall--Young calibration with the same pathways and settings. Because MAGMA gene--gene correlations depend only on genotypes, one correlation file was used for all permutations. Since every pathway is null in every permutation, calibrated $p$-values should be uniform, and a family-wise $p \le 0.05$ should occur in at most $5\%$ of permutations.

\textbf{Sorghum stem composition.} To assess whether the gene score captures new biology (not just calibration), we analysed three stem-composition traits in the same panel \citep{Brenton2016}: Brix at physiological maturity (stem-juice soluble solids; $\nBapNBrix$ lines) and two dry-matter fractions, neutral detergent fibre (NDF; cellulose, hemicellulose and lignin; $\nBapNNDF$ lines) and non-fibre carbohydrate (NFC; mainly sugars and starch; $\nBapNNFC$ lines). Each trait was tested with GEMMA’s univariate linear mixed model \citep{Zhou2012GEMMA} (Wald test) via vcf2gwas v0.8.9 \citep{Vogt2022vcf2gwas} on the panel genotypes ($\nBapNlines$ lines), using the centred relatedness matrix (all chromosomes), the first three genotype PCs as covariates, and variants with minor allele frequency $\ge0.05$. Gene analysis, pathways, Westfall--Young calibration (self-contained and competitive), genomic control and the gene score matched the dry-biomass analysis: MAGMA SNP-wise mean per chromosome with $25\,\mathrm{kb}$ windows and the same correlation structure, SorghumBicolorCyc pathways (5--200 genes), $B=10{,}000$, and the same seed (\texttt{hpc/bap\_gwas/01}--\texttt{04}). Genes were ranked by gene score per trait, and we identified genes in the top 100 for all three traits. Ranks are reported among all genes and, as for dry biomass, among genes with evidence in all three layers.

\emph{The chromosome~4 locus.} For the shared locus (Section~3.7), the lead variant for each trait was the strongest signal between $63.85$ and $64.10\,\mathrm{Mb}$. Linkage disequilibrium ($r^2$) between lead variants was computed from panel genotypes using PLINK~1.9 \citep{Chang2015PLINK}. Allele effects at the Brix lead variant were taken from each trait’s mixed model and reported in trait SD units. To test whether the allele’s effect on fibre is explained by the sugar–fibre trade-off, NDF was regressed on allele count across lines by least squares, first alone (coefficient $b$) and then with Brix as a covariate ($b_{\mathrm{adj}}$). If the allele changed fibre only by changing sugar, $b_{\mathrm{adj}}$ would be near zero; if it acted on fibre independently of sugar, $b_{\mathrm{adj}}$ would equal $b$. The share of the fibre effect explained by the trade-off is $1-b_{\mathrm{adj}}/b$, the difference method of mediation analysis \citep{VanderWeele2016Mediation}; we report the corresponding difference between homozygotes ($2b$ and $2b_{\mathrm{adj}}$). Because Brix is measured with error and is itself changed by the allele, this share is approximate: it measures how much of the fibre effect the sugar change predicts, not a causal split. The regressions are descriptive and ignore relatedness; association evidence comes from the mixed model. Code: \texttt{hpc/bap\_gwas/05\_chr4\_locus\_LSF.sh}, \texttt{43\_bap\_chr4\_locus.R}, \texttt{44\_bap\_chr4\_numbers.R}.

\emph{Leave one gene out.} To measure how much of each candidate's rank comes
from its own association passing through its pathways, the leave-one-gene-out
check of Section~2.5 was applied to every gene with pathway evidence among the
top $100$ of each trait's gene score ($\nBapLogoNBrix$, $\nBapLogoNNDF$ and
$\nBapLogoNNFC$ genes for Brix, NDF and NFC), which include \textit{SbALDH3E1} and the thioredoxin reductase. For
each such gene and each tested pathway containing it, the pathway was scored
with and without the gene in one Westfall--Young run with the settings of the
trait's pathway analysis (gene-level genomic control, the trait's gene--gene
correlations, SorghumBicolorCyc, $B=10{,}000$, the same truncation grid and
seed; reduced sets of at least $4$ genes). Each gene's pathway evidence, pathway
score, gene score and rank were then recomputed with every other gene unchanged,
both among all genes and among the genes with evidence in all three layers. We
report each candidate's rank without the pathway layer, with it, and without
its own contribution, and, for all checked genes of a trait, the number that
stay in the top $100$ and their median ranks
(\texttt{hpc/bap\_gwas/06\_logo\_LSF.sh}, \texttt{47\_bap\_logo.R}).

\textbf{Gene-level genomic control.} For all real analyses (ten lipid trait--ancestry sets, sorghum dry biomass, and three stem-composition traits), we centred gene $p$-values before pathway testing. We converted each gene $p$-value to a $z$-score $z_g=\Phi^{-1}(1-p_g)$, computed from $\log p_g$ to preserve ordering below $10^{-300}$ (Section~2.4), subtracted the GWAS-wide median $\hat\delta$, and converted back: $p^{\mathrm{GC}}_g=1-\Phi(z_g-\hat\delta)$. This location genomic control \citep{Devlin1999} uses one constant per GWAS, preserves gene order, gene--gene correlations, and the null draws of Section~2.3, and sets the median gene $p$-value to $0.5$. We included it because permutation nulls showed GWAS-specific shifts in gene $z$-scores that a self-contained test can misread as pathway signal, while the correlated null lacks this shift (Section~3.6). If no shift is present, centring has no effect; if the GWAS is saturated, centring removes some true polygenic signal and effectively asks whether pathway genes are more associated than the genome-wide median gene (Section~4.2). The estimated shifts were $0.15$--$0.33$ in the eight non-European lipid sets, $1.38$ (LDL-C) and $1.77$ (HDL-C) in the European sets, $0.13$ for sorghum dry biomass, and $\nBapGCBrix$, $\nBapGCNDF$ and $\nBapGCNFC$ for Brix, NDF and NFC.

\textbf{Competitive test and the gene score.} For each real analysis, we ran the competitive test (Section~2.3) alongside the self-contained test on the same centred gene $p$-values and pathways, using $B=10{,}000$ circular shifts (\texttt{wy\_gc\_real.R} with \texttt{COMP=1}). To assess the impact on candidate genes, we recomputed the gene score (Section~2.5) with the pathway layer ranked by the competitive standardised statistic $z_k^{\mathrm{comp}}$ or by the joint statistic $\max(z_k, z_k^{\mathrm{comp}})$ (the weaker of the two) instead of $z_k$, keeping all other layers, weights, and pathway sets unchanged (\texttt{36\_prioritize\_comp.R}). We then compared positive-control recovery across the three versions (\texttt{37\_prioritize\_comp\_recovery.R}).

\textbf{Comparison with PoPS.} We compared our gene score with PoPS v0.2 \citep{Weeks2023PoPS} using European LDL-C and HDL-C data. PoPS was run with default settings (ridge regression on public gene features; MAGMA covariates and error covariance; feature selection at $p<0.05$; HLA removed) on MAGMA gene statistics computed as above (same reference panel, $\pm10$~kb window, SNP-wise mean model) but using the required Ensembl annotation ($18{,}383$ protein-coding genes). Known genes were Open Targets Genetics gold-standard genes \citep{Mountjoy2021L2G} from the 2019 and 2023 releases (LDL-C: LDL cholesterol measurement and hypercholesterolaemia; HDL-C: HDL cholesterol measurement): $15$ for LDL-C and $18$ for HDL-C, curated independently of this study. For genes scored by both methods ($17{,}736$; CATFISH genes matched to Ensembl by unique symbol), we ranked genes by PoPS, CATFISH, best SNP $p$-value, and MAGMA gene $p$-value, and counted known genes in the top $10$, $20$, $50$, $100$, and $200$. We also evaluated single loci: for each gold-standard lead variant, candidates were genes with a transcription start site within $500$~kb; a locus was “solved” if the top-ranked candidate was the gold-standard gene (\texttt{40\_gold\_standard\_lipids.py}, \texttt{41\_benchmark\_genes.R}, \texttt{42\_benchmark\_topk.R}).

\textbf{Benchmark across nine traits.} To evaluate the gene score beyond lipids, we added seven European GWAS with Open Targets gold-standard genes \citep{Mountjoy2021L2G} and GWAS Catalog harmonised summary statistics: coronary artery disease \citep{vanderHarst2018CAD}, breast cancer \citep{Michailidou2017BRCA}, Crohn's disease and ulcerative colitis \citep{deLange2017IBD}, atrial fibrillation \citep{Nielsen2018AF}, systemic lupus erythematosus \citep{Bentham2015SLE}, and UK Biobank asthma \citep{Valette2021Asthma}. Variants were limited to 1000 Genomes European rsIDs; per-variant sample sizes were used when available, otherwise the study total. MAGMA, CATFISH (self-contained test with gene-level genomic control; Reactome), and the gene score were run as for European lipid traits, and PoPS as above (\texttt{hpc/benchmark/01}--\texttt{04}, \texttt{hpc/pops}). We excluded type 2 diabetes (only a subset of the large European meta-analysis variants are in the GWAS Catalog) and schizophrenia and rheumatoid arthritis (no European genome-wide file). Known genes were each trait’s Open Targets gold-standard genes from the 2019 and 2023 releases (LDL-C and HDL-C as above). We compared six rankings: CATFISH; CATFISH without the pathway layer (pathway weight $0$; GWAS+MAGMA); best variant; MAGMA gene $p$-value; PoPS; and the mean CATFISH/PoPS rank. Rankings were computed within genes scored by both CATFISH and PoPS, and we counted known genes in the top $10$--$200$. The main analysis excluded the extended MHC (chr6, $25$--$34$~Mb, GRCh37) because PoPS excludes HLA by default and long-range LD dominates immune-trait rankings. We also report results for all genes and for genes in a tested pathway (\texttt{45\_benchmark\_topk\_all.R}).

\textbf{Top 50 compared with the strongest variant.} For each of the nine traits, the genome-wide top 50 of the gene score was compared with the top 50 genes ranked by their strongest variant (the GWAS layer alone), after removing the extended MHC and re-ranking the remaining genes. Genes in the first list but not the second were counted as lifted, and genes in the second but not the first as dropped. A lifted gene was attributed to the pathway layer when it was also outside the top 50 of the score without the pathway layer. We counted the gold-standard genes in each group (\texttt{50\_lifted\_demoted.py}).

\textbf{Literature review of the top-ranked genes.} For the ten highest-ranked genes of each of the nine traits (extended MHC removed), we searched the literature for evidence linking the gene to the trait and classified each as: a gene for which published genetic or functional evidence (Mendelian disease, coding or loss-of-function variants, or risk variants shown to change its expression or splicing) points to the gene itself; a gene at a locus reported for the trait whose causal gene is unresolved; or a gene at a locus where published evidence points to a neighbouring gene. Each classification and its references are in \texttt{49\_candidate\_refs.csv}, which \texttt{49\_top10\_candidates.py} merges into Supplementary Table~\ref{tab:top10}.

